% from ega4-2-fr.tex
\oldpage[IV-2]{106}

\begin{remarks}[5.7.8]
\label{IV.2.s5.7.8-fr}
Soit $X$ un préschéma localement noethérien de dimension 1; il revient alors
au même de dire que $X$ est un préschéma de Cohen-Macaulay, ou qu'il vérifie
$(S_1)$, ou qu'il vérifie une des propriétés $(S_n)$ pour $n\geq1$, en vertu
des définitions (5.7.1) et (5.7.2). En vertu de (5.7.5), il revient donc
encore au même, pour les préschémas de dimension 1, de dire que $X$ est un
préschéma de Cohen-Macaulay ou qu'il n'a pas de cycle premier associé
immergé. Par exemple, un préschéma localement noethérien réduit de dimension
1 est un préschéma de Cohen-Macaulay.
\end{remarks}

\begin{proposition}[5.7.9]
\label{IV.2.s5.7.9-fr}
Soient $A$, $B$ deux anneaux noethériens, $\rho:A\to B$ un homomorphisme
d'anneaux, $M$ un $B$-module tel que $M_{[\rho]}$ soit un $A$-module de type
fini. Soit $\mathfrak p$ un idéal premier de $A$; les idéaux premiers de $B$
au-dessus de $\mathfrak p$ et appartenant à $\operatorname{Supp}(M)$ sont en
nombre fini, et si $(\mathfrak q_i)_{1\leq i\leq n}$ est la famille de ces
idéaux, on a
\begin{align}
\dim_{A_{\mathfrak p}}(M_{\mathfrak p})
&=\sup_i\dim_{B_{\mathfrak q_i}}(M_{\mathfrak q_i}),
\tag{5.7.9.1}\label{IV.2.s5.7.9.1-fr}\\
\operatorname{prof}_{A_{\mathfrak p}}(M_{\mathfrak p})
&=\inf_i\operatorname{prof}_{B_{\mathfrak q_i}}(M_{\mathfrak q_i}).
\tag{5.7.9.2}\label{IV.2.s5.7.9.2-fr}
\end{align}
\end{proposition}

Si $\mathfrak b$ est l'annulateur de $M$, $B/\mathfrak b$ s'identifie à un
sous-$A$-module de $\operatorname{End}_A(M_{[\rho]})$, donc est de type fini
puisque $A$ est noethérien; il n'y a donc qu'un nombre fini d'idéaux premiers
de $B$ au-dessus de $\mathfrak p$ et contenant $\mathfrak b$, et ce sont
précisément ceux qui appartiennent à $\operatorname{Supp}(M)$. Remplaçant
$B$ par $B/\mathfrak b$, ce qui ne change pas les seconds membres de
(5.7.9.1) et (5.7.9.2) (par (0, 16.1.9) et (0, 16.4.8)), on peut donc
supposer que $B$ est une $A$-algèbre finie. Posons $S=A-\mathfrak p$, et
$B'=S^{-1}B$; $B'$ est un anneau semi-local noethérien dont les idéaux
maximaux sont $S^{-1}\mathfrak q_i$ ($1\leq i\leq n$) et comme
$M_{\mathfrak p}$ est un $A_{\mathfrak p}$-module de type fini, on a
$\dim_{A_{\mathfrak p}}(M_{\mathfrak p})=\dim_{B'}(M_{\mathfrak p})$
(0, 16.1.9). Cela étant, la relation (5.7.9.1) devient un cas particulier de
(0, 16.1.7.4). Pour démontrer la relation (5.7.9.2), on se ramène aussitôt
comme dans (0, 16.4.8) au cas où
$\operatorname{prof}_{A_{\mathfrak p}}(M_{\mathfrak p})=0$, et le même
raisonnement que dans (0, 16.4.8) montre que $M_{\mathfrak p}$ contient un
sous-$B'$-module de longueur finie, et par suite aussi un sous-$B'$-module
simple; mais un tel sous-module est nécessairement isomorphe au corps
résiduel de l'un des $B_{\mathfrak q_i}$, donc il y a au moins un indice $i$
tel que
$\operatorname{prof}_{B_{\mathfrak q_i}}(M_{\mathfrak q_i})=0$
(0, 16.4.6), ce qui termine la démonstration.

\begin{corollary}[5.7.10]
\label{IV.2.s5.7.10-fr}
Supposons vérifiées les hypothèses de (5.7.9), et supposons en outre que $A$
soit un anneau local; alors, pour que $M$ soit un $A$-module de
Cohen-Macaulay, il faut et il suffit que, pour tous les idéaux premiers
$\mathfrak q_i$ de $B$ au-dessus de l'idéal maximal de $A$,
$M_{\mathfrak q_i}$ soit un $B_{\mathfrak q_i}$-module de Cohen-Macaulay, et
de plus que tous les nombres
$\dim_{B_{\mathfrak q_i}}(M_{\mathfrak q_i})$ soient égaux.
\end{corollary}

Il résulte en effet de (5.7.9.1) et de (5.7.9.2) que ces conditions
équivalent à la relation
$\dim_A(M)=\operatorname{prof}_A(M)$.

\begin{corollary}[5.7.11]
\label{IV.2.s5.7.11-fr}
Supposons vérifiées les hypothèses de (5.7.9).
\begin{enumerate}
\item[(i)] Si $M_{[\rho]}$ vérifie la propriété $(S_k)$, il en est de même
de $M$.
\item[(ii)] Supposons que pour tout couple d'idéaux premiers
$\mathfrak q$, $\mathfrak q'$ de $B$ tels que
$\rho^{-1}(\mathfrak q)=\rho^{-1}(\mathfrak q')$, on ait
$\dim_{B_{\mathfrak q}}(M_{\mathfrak q})
=\dim_{B_{\mathfrak q'}}(M_{\mathfrak q'})$. Alors, si $M$ vérifie la
propriété $(S_k)$, il en est de même de $M_{[\rho]}$.
\end{enumerate}
\end{corollary}

